\documentclass[11pt]{article}
\usepackage[margin=1in]{geometry}
\usepackage{amsmath,amssymb,amsthm}
\usepackage{hyperref}

\newtheorem{theorem}{Theorem}
\newtheorem{proposition}{Proposition}

\newcommand{\K}{\mathbb K}
\newcommand{\R}{\mathbb R}
\newcommand{\calP}{\mathcal P}
\newcommand{\calB}{\mathcal B}

\title{Ordinary Pietsch measures already characterize injective dominated polynomials}
\author{fa\_banach\_001, agent\_lane\_14}
\date{2026-06-29}

\begin{document}
\maketitle

\section*{Source and Question}

The source paper is Botelho--Pellegrino--Rueda, \emph{On Pietsch
measures for summing operators and dominated polynomials},
arXiv:1210.3332.  Section 3 asks for a polynomial analogue of the
linear injective Pietsch-measure criterion:

\begin{quote}
Given \(n\in\mathbb N\), is it true that \(j_p^e\) is injective if and
only if \(\mu\) is a Pietsch measure for some injective
\(p\)-dominated \(n\)-homogeneous polynomial defined on \(E\)?
\end{quote}

The source then observes that, for \(n\ge2\), injective
\(n\)-homogeneous polynomials can occur only in the real odd-degree
case, and proves an iff statement with the stronger hypothesis that
\(\mu\) is a \emph{differential} Pietsch measure.  The result below
answers the stated question in the ordinary Pietsch-measure form in the
non-vacuous normed case.

\section*{Notation}

Let \(E\) be a Banach space over \(\K\).  The dual unit ball
\(B_{E'}\) is endowed with the weak-star topology, and
\[
  e:E\longrightarrow C(B_{E'}),\qquad e(x)(\varphi)=\varphi(x),
\]
is the canonical isometry.  If \(\mu\) is a regular Borel probability
measure on \(B_{E'}\), let
\[
  j_p:C(B_{E'})\longrightarrow L_p(\mu)
\]
be the canonical quotient map and write \(j_p^e\) for its restriction
to \(e(E)\).

For an \(n\)-homogeneous polynomial \(P:E\to F\), saying that \(\mu\)
is a Pietsch measure for \(P\) means that there is a constant \(C\)
such that
\[
  \|P(x)\|\le C
  \left(\int_{B_{E'}} |\varphi(x)|^p\,d\mu(\varphi)\right)^{n/p}
  \qquad (x\in E).
\]

\section*{Proof Intuition}

The source uses differential Pietsch measures in order to prove the
stronger statement that \(j_p^e(e(x))=j_p^e(e(y))\) implies
\(P(x)=P(y)\).  For the injectivity criterion this is more than is
needed.  If \(j_p^e(e(x))=j_p^e(e(y))\), then \(x-y\) has zero Pietsch
seminorm.  Ordinary domination applied to \(x-y\) gives \(P(x-y)=0\).
Since a homogeneous polynomial satisfies \(P(0)=0\), injectivity of
\(P\) forces \(x-y=0\).

The converse direction is the canonical odd-power construction already
implicit in the source: if \(j_p^e\) is injective, then
\[
  x\longmapsto (j_{p/n}^e(e(x)))^n
\]
is an injective \(n\)-homogeneous polynomial into \(L_{p/n}(\mu)\) when
the scalars are real and \(n\) is odd.

\begin{proposition}
\label{prop:ordinary-implies-injective}
Let \(E\) be a Banach space over \(\K\), let \(0<p<\infty\), let
\(\mu\) be a regular Borel probability measure on \(B_{E'}\), and let
\(P\in\calP(^nE;F)\) be an injective \(n\)-homogeneous polynomial.  If
\(\mu\) is a Pietsch measure for \(P\), then \(j_p^e:e(E)\to L_p(\mu)\)
is injective.
\end{proposition}

\begin{proof}
Suppose \(j_p^e(e(x))=j_p^e(e(y))\).  Since \(e\) is linear, this says
that \(e(x-y)\) is zero in \(L_p(\mu)\), hence
\[
  \int_{B_{E'}}|\varphi(x-y)|^p\,d\mu(\varphi)=0.
\]
Applying the Pietsch domination inequality for \(P\) to \(x-y\) gives
\[
  \|P(x-y)\|\le
  C\left(\int_{B_{E'}}|\varphi(x-y)|^p\,d\mu(\varphi)\right)^{n/p}=0.
\]
Thus \(P(x-y)=0=P(0)\).  Since \(P\) is injective, \(x-y=0\), so
\(x=y\).  Therefore \(j_p^e\) is injective.
\end{proof}

\begin{theorem}
\label{thm:ordinary-criterion}
Let \(E\) be a real Banach space, let \(n\in\mathbb N\) be odd, let
\(p\ge n\), and let \(\mu\) be a regular Borel probability measure on
\(B_{E'}\).  Then \(j_p^e\) is injective if and only if \(\mu\) is a
Pietsch measure for some injective \(p\)-dominated \(n\)-homogeneous
polynomial defined on \(E\).
\end{theorem}

\begin{proof}
The reverse implication is Proposition~\ref{prop:ordinary-implies-injective}.

Conversely, assume that \(j_p^e\) is injective.  Since \(p\ge n\), the
space \(L_{p/n}(\mu)\) is a Banach space.  Define
\[
  P:E\longrightarrow L_{p/n}(\mu),\qquad
  P(x)=\big(j_{p/n}(e(x))\big)^n.
\]
This is a continuous \(n\)-homogeneous polynomial.  Moreover,
\[
  \|P(x)\|_{L_{p/n}(\mu)}
  =
  \left(\int_{B_{E'}}|\varphi(x)|^p\,d\mu(\varphi)\right)^{n/p},
\]
so \(\mu\) is a Pietsch measure for \(P\) with constant \(1\).  By the
Pietsch domination theorem for dominated homogeneous polynomials, \(P\)
is \(p\)-dominated.

It remains to check injectivity.  If \(P(x)=P(y)\), then
\[
  (j_{p/n}(e(x)))^n=(j_{p/n}(e(y)))^n
  \quad\text{in }L_{p/n}(\mu).
\]
Over the real scalars with \(n\) odd, the scalar map \(t\mapsto t^n\)
is injective, so \(j_p(e(x))=j_p(e(y))\) in \(L_p(\mu)\).  Since
\(j_p^e\) is injective, \(x=y\).  Hence \(P\) is injective.
\end{proof}

\section*{Limitations and Novelty Check}

The theorem is stated for real odd \(n\) and \(p\ge n\), matching the
non-vacuous normed setting isolated in the source.  For complex
scalars and \(n\ge2\), and for real even \(n\), an injective
\(n\)-homogeneous polynomial cannot exist because nontrivial roots of
unity (or the sign change \(x\mapsto -x\)) identify distinct points.

Cheap run indexes were searched for arXiv:1210.3332, Pietsch measures,
differential Pietsch measures, \(j_p^e\), and injective dominated
polynomials.  Web searches for exact title and close phrases found the
source arXiv record and no later paper with this ordinary-measure
refinement.  The result should be reviewed as a small strengthening of
the source's Proposition 3.4: the differential Pietsch-measure
hypothesis is unnecessary for the injective-polynomial criterion.

\begin{thebibliography}{9}
\bibitem{BPR}
G. Botelho, D. Pellegrino, and P. Rueda,
\emph{On Pietsch measures for summing operators and dominated
polynomials}, arXiv:1210.3332.
\end{thebibliography}

\end{document}
